% Plantilla mínima para TFG (TFG / Trabajo Fin de Grado)
% Archivo: Info_TFG.tex
% Cómo usar: editar título/autor y dividir en capítulos con \include o \input si se desea.
% Compilar (recomendado): usar latexmk o pdflatex + biber/bibtex según bibliografía.
% Ejemplo PowerShell (TeX Live / MiKTeX instalado):
%   latexmk -pdf Info_TFG.tex
% o, paso a paso:
%   pdflatex Info_TFG.tex; biber Info_TFG; pdflatex Info_TFG.tex; pdflatex Info_TFG.tex

\documentclass[12pt,a4paper,oneside]{report}

% Codificación y idioma
\usepackage[T1]{fontenc}
\usepackage[utf8]{inputenc}
\usepackage[spanish]{babel}

% Márgenes y microtipografía
\usepackage[a4paper,margin=2.5cm]{geometry}
\usepackage{microtype}

% Gráficos, tablas y matemáticas
\usepackage{graphicx}
\usepackage{caption}
\usepackage{subcaption}
\usepackage{booktabs}
\usepackage{amsmath,amssymb}

% Hipervínculos
\usepackage[hidelinks]{hyperref}

% Numeración de secciones sin dependencia de capítulos (1,2,3... en vez de 0.1,0.2)
\usepackage{chngcntr}
\counterwithout{section}{chapter}

% Configuración del título/portada (personaliza según la normativa de tu universidad)
\title{Comparación de Nariz Electrónica con olfato de Agentes Caninos para la detección de Drogas}
\author{Raúl Martín-Romo Sánchez\\
	\textit{Grado en Ingeniería Informática del Software}\\
	\textbf{Tutor:} Pablo García Rodríguez\\
    \textbf{Co-tutor:} Jesús Lozano Rogado}
\date{\today}

\begin{document}

% Portada simple
\maketitle

% Capítulos de ejemplo
\chapter{Idea del proyecto}
\label{ch:idea_proyecto}
Desarrollar una nariz electrónica capaz de detectar un tipo de droga ilegal y compararla con el olfato de los agentes caninos entrenados del cuerpo de detección de drogas de la Guardia Civil de Cáceres.

\chapter{Motivación}
\label{ch:motivacion}
Revisión bibliográfica y trabajos relacionados.

\chapter{Objetivos}
\label{ch:objetivos}
Describe los métodos, herramientas y procedimientos empleados.

\chapter{Materiales}
\label{ch:materiales}
Sensores MOX (Sensores de óxidos metálicos) y materiales por identificar tras la formación con Jesús Lozano Rogado.

\chapter{Planificación}
\label{ch:planificacion}
Describe la planificación temporal y las fases del proyecto.

\chapter{Resultados}
\label{ch:resultados}
Presenta resultados, tablas, figuras y su discusión.

\end{document}

